\documentclass[10pt,a4paper]{amsart}
\usepackage[margin=19mm]{geometry}
\usepackage{amsmath,amssymb,mathtools}
\usepackage[scr=rsfs,cal=euler]{mathalfa}
\usepackage{xfrac}
\usepackage[unicode,hidelinks,bookmarks=false]{hyperref}
\newcommand{\ie}{i.e.\ }
\newcommand{\cf}{cf.\ }
\newcommand{\resp}{respectively\ }
\newcommand{\Subsection}{\subsection}
\providecommand{\nobreakdash}{\nobreak\hbox{-}}
\setlength{\emergencystretch}{2em}
\multlinegap0pt
\usepackage{amssymb}
\usepackage{mathtools}
\usepackage{enumerate}
\newtheorem{thm}{Theorem}
\newcommand{\B}{\mathcal{B}}
\newcommand{\NN}{\mathbb{N}}
\DeclareMathOperator*{\R}{Re}
\newcommand{\RR}{\mathbb{R}}
\newcommand{\dd}{\mathrm{d}}
\newcommand{\ep}{\varepsilon}
\newcommand{\inv}{^{-1}}
\title{Polynomial stability of non-linearly damped contraction semigroups}
\author{Lassi Paunonen, David Seifert}
\date{Source: arXiv:2509.04275v2 (CC BY 4.0)}
\begin{document}
\maketitle

\noindent\emph{Source and licence.} Polynomial stability of non-linearly damped contraction semigroups. Version arXiv:2509.04275v2 (CC BY 4.0), distributed by arXiv under the Creative Commons Attribution 4.0 International licence, http://creativecommons.org/licenses/by/4.0/, as recorded in the arXiv licence metadata for this exact version and shown on the abstract page https://arxiv.org/abs/2509.04275v2. Full licence notice: \texttt{licenses/arxiv-2509.04275-CC-BY-4.0.txt}.\par\smallskip

\noindent\textit{Editorial scope.} Counted unit: the article's thm thm:main (source0.tex lines 145--167). The article's complete original proof of it is delivered with it (source0.tex lines 169--252, one \texttt{proof} environment of 84 lines). No mathematics is written, repaired or re-derived: the statement and the proof are the article's own lines, in the article's own notation, and no retained line is substituted or re-flowed.  Retained uncounted as the source-local context the proof needs: lines 82--84; lines 90--97; lines 121--124.  Nothing was omitted from any retained span, so the package carries no omission record.  The retained lines cite 9 works, whose entries are transcribed verbatim from the article's own bibliography source in the e-print and delivered with short editorial labels recorded in the item's reference mapping.  No retained line contains a figure environment, an image-inclusion command or drawing code, so no image is delivered and the package declares no asset dependency.  Editorial formatting only, in addition: an AMS \texttt{amsart} preamble that restates the article's own declarations and macros used by the retained lines, namely the environments \texttt{thm} on separate counters, together with the standard packages those lines need.  The retained environments print in this document as Theorem 1 in order of appearance, whereas the article prints the same items with the numbering of its own full document.  The complete native source of the article as distributed in the arXiv e-print for this exact version is bundled unchanged as \texttt{sources/184478/source0.tex}.
\begin{equation}\label{eq:sys}
\dot x(t)=Ax(t) - B\phi(B^\ast x(t)), \qquad t\ge0,
\end{equation}

\begin{equation}
\label{eq:PhiAssStab}
\R\langle\phi(u),u\rangle \gtrsim
\begin{cases}
 \|u\|^2, &   \|u\|\le\delta,\\
1,  &\|u\|>\delta,
\end{cases}
\end{equation}

\begin{equation}\label{eq:energy_balance}
\|x(t)\|^2 + 2\int_0^t \R \langle \phi(B^\ast x(s)),B^\ast x(s)\rangle\,\dd s 
\le \|x(0)\|^2, \qquad t\ge 0,
\end{equation}

\begin{thm}\label{thm:main}
Suppose that $\phi\colon U\to U$ is monotone, locally Lipschitz continuous and satisfies $\phi(0)=0$, and that~\eqref{eq:PhiAssStab}
holds for some $\delta>0$.
Suppose further that there exist $\beta,\tau,c_\tau>0$ such that
\begin{equation}
\label{eq:ObsEstimate}
c_\tau\|(I-A)^{-\beta} x_0\|^2 
\le \int_0^\tau \|B^\ast T(t)x_0\|^2 \,\dd t, \qquad x_0\in X.
\end{equation}
Then
 $i\RR\subseteq\rho(A_B)$. Furthermore,
 all  mild solutions of~\eqref{eq:sys} satisfy $\|x(t)\|\to 0$ as $t\to\infty$, and
 all classical solutions of~\eqref{eq:sys} satisfy $\|x(t)\|=O(t^{-1/(2\beta)})$ as $t\to\infty$.

If, in addition,
$\|(is-A_B)\inv\|\lesssim 1+|s|^\alpha$ for some $\alpha>0$ and all $s\in\RR$, and
 if there exist $\gamma>\alpha/2+1$ and $\kappa,\ep>0$ such that
\begin{equation}
\label{eq:PhiApproximatLinearity}
\|\phi(u)-\kappa u\|\lesssim \|u\|^\gamma, \qquad  \|u\|\le \ep ,
\end{equation}
then $\|x(t)\|= O(t^{-1/\alpha})$ as $t\to\infty$ for any classical solution of~\eqref{eq:sys}.
\end{thm}

\begin{proof}
Let $x(\cdot)$ be a classical solution of~\eqref{eq:sys}, and suppose there exists $t_0\ge0$ such that $\|B^*x(t)\|\le\delta$ for all $t\ge t_0$. Furthermore, let $L_\delta>0$ be such that $\|\phi(u)\|\le L_\delta\|u\|$ for all $u\in U$ with $\|u\|\le\delta$. Let $\tau>0$. For $u\in L^2(0,\tau;U)$ we define 
 $$(\mathbb{F}_\tau u)(t)=\int_0^t B^\ast T(t-s)Bu(s)\,\dd s,\qquad 0\le t\le\tau.$$
Then $\mathbb{F}_\tau\in \B(L^2(0,\tau;U))$ and 
$$B^\ast T(s)x(t) = B^\ast x(t+s)+ (\mathbb{F}_\tau\phi(B^\ast x(t+\cdot)))(s),\qquad 0\le s\le\tau,\ t\ge0,$$
by the variation of parameters formula, and hence
$$
\|B^\ast T(\cdot)x(t)\|_{L^2(0,\tau;U)}\lesssim\|B^\ast x(t+\cdot)\|_{L^2(0,\tau;U)}+\|\phi(B^\ast x(t+\cdot))\|_{L^2(0,\tau;U)}
$$
for all $t\ge0$. For $t\ge t_0$ it follows that 
$$\|B^\ast T(\cdot)x(t)\|_{L^2(0,\tau;U)}\lesssim\|B^\ast x(t+\cdot)\|_{L^2(0,\tau;U)},$$
 and combining this with~\eqref{eq:ObsEstimate} gives
$$\|(I-A)^{-\beta} x(t)\|
\lesssim  \|B^\ast x(t+\cdot)\|_{L^2(0,\tau;U)}, \qquad t\ge t_0.$$
Now let $k\in\NN$ be such that $k\tau\ge t_0$. Using~\eqref{eq:energy_balance} and~\eqref{eq:PhiAssStab} it follows that
$$\begin{aligned}
\|x(k\tau)\|^2&-\|x((k+1)\tau)\|^2= 2\int_0^\tau \R \langle \phi(B^\ast x(k\tau+s)),B^\ast x(k\tau+s)\rangle\, \dd s\\
&\gtrsim \|B^\ast x(k\tau+\cdot))\|_{L^2(0,\tau;U)}^2\gtrsim\|(I-A)^{-\beta}x(k\tau) \|^2,
\end{aligned}
$$
where the implicit constants are independent of $k$. Noting that $x(k\tau)\in D(A)$, we may apply the moment inequality~\cite[Thm.~II.5.34]{EngNag00book} to obtain
$$\|x(k\tau)\|\le \|(I-A)x(k\tau)\|^{\frac{\beta}{1+\beta}}\|(I-A)^{-\beta}x(k\tau)\|^{\frac{1}{1+\beta}}.$$
Thus
$$
\|x(k\tau)\|^2-\|x((k+1)\tau)\|^2\gtrsim\frac{\|x((k+1)\tau)\|^{2(1+\beta)}}{\|(I-A)x(k\tau)\|^{2\beta}},
$$
where we have used the fact that $\|x((k+1)\tau)\|\le\|x(k\tau)\|$. Now
$$\begin{aligned}
\|(I-A)x(k\tau)\|&\le \|x(k\tau)\|+\|Ax(k\tau)-B\phi(B^* x(k\tau))\|+\|B\|\|\phi(B^*x(k\tau))\|\\
&\le (1+L_\delta\|B\|^2)\|x(k\tau)\|+\|\dot{x}(k\tau)\|\\
&\le (1+L_\delta\|B\|^2)\|x(t_0)\|+\|\dot{x}(t_0)\|\\
&=(1+L_\delta\|B\|^2)\|x(t_0)\|+\|Ax(t_0)-B\phi(B^*x(t_0))\|\\
&\le 2(1+L_\delta\|B\|^2)\|x(t_0)\|+\|(I-A)x(t_0)\|\\
&\lesssim\|(I-A)x(t_0)\|,
\end{aligned}
$$
where we have used monotonicity of the functions $\|x(\cdot)\|$ and $\|\dot{x}(\cdot)\|$. It follows that
$$\frac{\|x((k+1)\tau)\|^2}{\|(I-A)x(t_0)\|^2}\le\frac{\|x(k\tau)\|^2}{\|(I-A)x(t_0)\|^2}-c\left(\frac{\|x((k+1)\tau)\|^2}{\|(I-A)x(t_0)\|^2}\right)^{1+\beta}$$
for some constant $c>0$ and all sufficiently large $k\ge1$, and hence
$$\|x(k\tau)\|\lesssim \frac{\|(I-A)x(t_0)\|}{(k\tau+1)^{1/(2\beta)}}$$
for all sufficiently large $k\ge1$ by~\cite[Lem.~1.3.4]{AmmNic15book}. By monotonicity of $\|x(\cdot)\|$ we deduce that 
$$\|x(t)\|\lesssim \frac{\|(I-A)x(t_0)\|}{(t+1)^{1/(2\beta)}}$$
for all sufficiently large $t\ge0$. Next we observe that, for any $\delta\ge\|B^*\|\|x_0\|$, the identity function on $U$ defined by $\phi(u)=u$ for all $u\in U$ satisfies our assumptions  with $t_0=0$ and $L_\delta=1$. We deduce that
$\|T_B(t)(I-A_B)\inv\|=O(t^{-1/(2\beta)})$ as $t\to\infty$, where $(T_B(t))_{t\ge0}$ is the $C_0$-semigroup of contractions generated by $A_B$. It follows from~\cite[Thm.~1.1]{BatDuy08} that $i\RR\subseteq\rho(A_B)$. Moreover, $(T_B(t))_{t\ge0}$ is strongly stable in the sense that $\|T_B(t)x\|\to0$ as $t\to\infty$ for all $x\in X$, by a standard density argument. It then follows from~\cite[Thm.~2.2]{CurZwa16} (noting that the result carries over, with the appropriate modifications, to the setting of complex Hilbert spaces, and that the assumptions of compact resolvent and approximate observability can be replaced by strong stability of $(T_B(t))_{t\ge0}$) that whenever $\phi$ satisfies the conditions of our theorem we have $\|x(t)\|\to0$ as $t\to\infty$ for all  mild solutions of~\eqref{eq:sys}. Hence the first part of the proof shows that $\|x(t)\|=O(t^{-1/(2\beta)})$ as $t\to\infty$ for  all classical solutions of~\eqref{eq:sys}, which completes the proof of the first part of the result.

Now assume, in addition, that
$\|(is-A_B)\inv\|\lesssim 1+|s|^\alpha$ for some $\alpha>0$ and all $s\in\RR$, and that
 there exist $\gamma>\alpha/2+1$ and $\kappa,\ep>0$ such that~\eqref{eq:PhiApproximatLinearity} holds. If $\alpha\ge 2\beta$ then the result already follows without any of the additional assumptions from what has already been proved, so we may assume that $\alpha<2\beta$. Recall from the first part of the proof that $\|x(t)\|\to0$ as $t\to\infty$ for all  mild solutions of~\eqref{eq:sys}, and that $\|x(t)\|=O(t^{-1/(2\beta)})$ as $t\to\infty$ for all classical solutions of~\eqref{eq:sys}. Note also that by~\cite[Lem.~2.11(c)]{ChiPau23} the operator $A_\kappa = A-\kappa BB^\ast$ satisfies $i\RR\subseteq\rho(A_\kappa)$ and $\|T_\kappa(t)A_\kappa\inv\|=O(t^{-1/\alpha})$ as $t\to\infty$, where $(T_\kappa(t))_{t\ge0}$ is the contraction semigroup  generated by $A_\kappa$. Let $x(\cdot)$ be a classical solution of~\eqref{eq:sys}, and let $t_0\ge0$ be such that $\|B^*x(t)\|\le\ep$ for all $t\ge t_0$.  We have
 $$\dot{x}(t)=(A-\kappa BB^\ast) x(t) + By(t),\qquad t\ge 0,$$
 where $y(t)=\kappa B^\ast x(t)-\phi(B^\ast x(t))$ for all $t\ge 0$. Note that, by~\eqref{eq:PhiApproximatLinearity} and boundedness of $B$, we have $\|y(t)\|\lesssim \| x(t)\|^{\gamma}$ for $t\ge t_0$. Given $t\ge0$ let $t_\theta=\max\{t_0,t-t^\theta/2\}$, where $\theta\in(0,1]$ is to be chosen later, noting that $t_\theta\to\infty$ and $t-t_\theta\to\infty$ as $t\to\infty$. The variation of parameters formula gives
 \begin{equation}\label{eq:three}
 \begin{aligned}
 x(t)
 = T_\kappa (t-t_0)x(t_0) 
&+ T_\kappa (t-t_\theta)\int_{t_0}^{t_\theta} T_\kappa (t_\theta-s)By(s)\,\dd s\\
\qquad &+ \int_{t_\theta}^t T_\kappa (t-s)By(s)\,\dd s
 \end{aligned}
 \end{equation}
for all $t\ge t_0$.  Since $x(t)\in D(A)$ for all $t\ge 0$, we have  $\|T_\kappa (t-t_0)x(t_0)\| =O(  t^{-1/\alpha})$ as $t\to\infty$. 
 Moreover, since $\phi$ is locally Lipschitz continuous and $Bx(\cdot)$ is continuously differentiable, we have $\int_{t_0}^{t_\theta} T_\kappa (t_\theta-s)By(s)\,\dd s\in D(A)$ for all $t\ge t_0$ by~\cite[Cor.~3.1.17]{AreBat11book}.
The second term in~\eqref{eq:three} therefore satisfies
$$\begin{aligned}
\bigg\|T_\kappa (t-t_\theta)\int_{t_0}^{t_\theta} T_\kappa (t_\theta&-s)By(s)\,\dd s\bigg\|\\ &\le\|T_\kappa (t-t_\theta)A_\kappa\inv\|\bigg\|A_\kappa\int_{t_0}^{t_\theta} T_\kappa (t_\theta-s)By(s)\,\dd s\bigg\|
\end{aligned}$$
 for all $t\ge t_\theta$. Now $\|T_\kappa (t-t_\theta)\smash{A_\kappa\inv}\|=O(t^{-\theta/\alpha})$ as $t\to\infty$. On the other hand,  we have 
$$ \int_{t_0}^{t_\theta} T_\kappa (t-s)By(s)\,\dd s = x(t_\theta)-T_\kappa (t_\theta-t_0)x(t_0),\qquad t\ge t_0,$$ 
by the variation of parameters formula. Thus 
$$\begin{aligned}
\bigg\|A_\kappa&\int_{t_0}^{t_\theta} T_\kappa (t_\theta-s)By(s)\,\dd s\bigg\|\le \|A_\kappa x(t_\theta)\|+\|A_\kappa T_\kappa (t_\theta-t_0)x(t_0)\|\\
&\le  \|Ax(t_\theta) - B\phi(B^\ast x(t_\theta))\|+\|By(t_\theta)\|+\|T_\kappa (t_\theta-t_0)A_\kappa x(t_0)\|\\
&\lesssim \|\dot{x}(t_0)\|+\|x(t_0)\|^\gamma+\|A_\kappa x(t_0)\|
\end{aligned}$$
for all $t\ge t_0$, where we have used contractivity of $(T_\kappa(t))_{t\ge0}$ and the fact that both $\|x(\cdot)\|$ and $\|\dot x(\cdot)\|$ are non-increasing functions. Thus the second term in~\eqref{eq:three} satisfies
$$\bigg\|T_\kappa (t-t_\theta)\int_{t_0}^{t_\theta} T_\kappa (t_\theta-s)By(s)\,\dd s\bigg\|=O({t^{-\theta/\alpha}}),\qquad t\to\infty.$$
The third and final term in~\eqref{eq:three} satisfies
$$ \int_{t_\theta}^t T_\kappa (t-s)By(s)\,\dd s
= \int_{0}^t T_\kappa (t-s)By_\theta(s)\,\dd s,$$
where $y_\theta(t)=0$ for $t\in[0,t_\theta)$ and $y_\theta(t)=y(t)$ for $t\ge t_\theta$. By~\cite[Lem.~2.2.6]{Oos00book} the maps $\Phi_t\in\B(L^2(0,t;U),X)$, defined by $\Phi_t u=\int_0^t T_\kappa (t-s)Bu(s)\,\dd s$ for $t\ge 0$ and $u\in L^2(0,t;U)$, have uniformly bounded operator norms; see also~\cite[Thm.~6.5.6]{CurZwa20book}, \cite[Cor.~6.1]{Sta02}. Since $\|y(t)\|\lesssim \|x(t)\|^\gamma$ for all $t\ge t_0$, and since $\|x(t)\|=O(t^{-1/(2\beta)})$ as $t\to\infty$ by the first part of the result, we deduce that
$$\bigg\| \int_{t_\theta}^t T_\kappa (t-s)By(s)\,\dd s\bigg\|^2\lesssim  \int_{t_\theta}^t \|y(s)\|^2\,\dd s\lesssim\int_{t_\theta}^t\frac{\dd s}{(1+s)^{\gamma/\beta}}\le\frac{t-t_\theta}{(1+t_\theta)^{\gamma/\beta}}$$
for all $t\ge t_0$.
It follows that the three terms in~\eqref{eq:three} are of order
$O(t^{-1/\alpha})$, $O(t^{-\theta/\alpha})$ and  $O(t^{-\mu})$, respectively, as $t\to\infty$, where  $\mu=\gamma/(2\beta)-\theta/2$. If  $\beta\le \alpha\gamma/(\alpha+2)$
we set $\theta=1$. Then $\mu\ge1/\alpha$, so the third term decays at least as fast as the first two, and we obtain $\|x(t)\|=O(t^{-1/\alpha})$ as $t\to\infty$, which is what we wanted to prove. On the other hand, if $\beta> \alpha\gamma/(\alpha+2)$ we set $\theta=\alpha\gamma/(\beta(\alpha+2))$. Then $\theta\in(0,1)$ and the rate is determined by the second and third terms, which decay at the same speed, giving $\|x(t)\|=O(t^{-1/(2\sigma\beta)})$ as $t\to\infty$, where $\sigma=(\alpha+2)/(2\gamma)$. The latter estimate is strictly worse than the one we wish to prove. On the other hand, since $\gamma>\alpha/2+1$ by assumption, we have $\sigma\in(0,1)$ and hence this decay rate is strictly \emph{better} than the estimate $\|x(t)\|=O(t^{-1/(2\beta)})$ as $t\to\infty$ coming from the first part of the result. In this second case, we may therefore repeat the argument with $\beta$ replaced by $\sigma\beta$. As above, we find that if $\sigma\beta\le \alpha\gamma/(\alpha+2)$ then $\|x(t)\|=O(t^{-1/\alpha})$ as $t\to\infty$ and the proof is complete, while if $\sigma\beta> \alpha\gamma/(\alpha+2)$ then  $\|x(t)\|=O\smash{(t^{-1/(2\sigma^2\beta)})}$ as $t\to\infty$. If necessary we may now iterate this process, terminating after $k\ge0$ repetitions if $\sigma^k\beta\le \alpha\gamma/(\alpha+2)$, giving the desired the rate $\|x(t)\|=O(t^{-1/\alpha})$ as $t\to\infty$, and otherwise replacing $\sigma^k\beta$ by $\sigma^{k+1}\beta$. Since $\sigma^k\beta\le \alpha\gamma/(\alpha+2)$ for all sufficiently large $k\ge0$ by virtue of the fact that $\sigma\in(0,1)$, the process must eventually terminate yielding $\|x(t)\|=O(t^{-1/\alpha})$ as $t\to\infty$, as required.
\end{proof}

\begin{thebibliography}{99}
\bibitem[AmmNic]{AmmNic15book} K.~Ammari and S.~Nicaise. \newblock {\em Stabilization of elastic systems by collocated feedback}, volume 2124 of {\em Lecture Notes in Mathematics}. \newblock Springer, Cham, 2015.
\bibitem[AreBat]{AreBat11book} W.~Arendt, C.J.K. Batty, M.~Hieber, and F.~Neubrander. \newblock {\em Vector-valued Laplace Transforms and Cauchy Problems}. \newblock Birkh{\"a}user, Basel, second edition, 2011.
\bibitem[BatDuy]{BatDuy08} C.J.K. Batty and T.~Duyckaerts. \newblock Non-uniform stability for bounded semi-groups on {B}anach spaces. \newblock {\em J. Evol. Equ.}, 8:765--780, 2008.
\bibitem[ChiPau]{ChiPau23} R.~{Chill}, L.~{Paunonen}, D.~{Seifert}, R.~{Stahn}, and Yu. {Tomilov}. \newblock Non-uniform stability of damped contraction semigroups. \newblock {\em Anal. PDE}, 16(5):1089--1132, 2023.
\bibitem[CurZwa]{CurZwa16} R.F. Curtain and H.~Zwart. \newblock Stabilization of collocated systems by nonlinear boundary control. \newblock {\em Systems Control Lett.}, 96:11--14, 2016.
\bibitem[CurZwa]{CurZwa20book} R.F. Curtain and H.~Zwart. \newblock {\em Introduction to Infinite-Dimensional Systems Theory}, volume~71 of {\em Texts in Applied Mathematics}. \newblock Springer, New York, 2020.
\bibitem[EngNag]{EngNag00book} K.-J. Engel and R.~Nagel. \newblock {\em One-Parameter Semigroups for Linear Evolution Equations}. \newblock Springer, New York, 2000.
\bibitem[Oos00b]{Oos00book} J.~Oostveen. \newblock {\em Strongly Stabilizable Distributed Parameter Systems}. \newblock Society for Industrial and Applied Mathematics (SIAM), Philadelphia, PA, 2000.
\bibitem[Sta02]{Sta02} O.~Staffans. \newblock Passive and conservative continuous-time impedance and scattering systems. {Part I}: {W}ell-posed systems. \newblock {\em Math. Control Signals Systems}, 15(4):291--315, 2002.
\end{thebibliography}
\end{document}
